% Short running form: the full title is wider than the box the running head
% leaves beside the Sunday's name.
\runninghead{The Lord Who Grasps the Hand of Kings}
\properlane{cyrus-and-caesar}{The Lord Who Grasps the Hand of Kings: Cyrus, Caesar and the Limit of the Debt}
\label{sec:cyrus-and-caesar}

The \work{Ordo} sets two titles side by side: the Lord grasps the right hand of
Cyrus to subdue nations before him, and Caesar is to be given what is
Caesar's. This interpretation reads the Gospel through that pairing, the one
relation between the readings that the liturgical books themselves declare.
Jerome and Aquinas explain why a gentile king who does not know the Lord is
called his anointed, and what the Lord required of his power; John Chrysostom says that tribute is
rendered as a debt and that the debt ends where godliness would be harmed;
Augustine shows Christian soldiers who served an unbelieving emperor and
refused him only in the cause of Christ; Aquinas says that the hearers are at
once God's and Caesar's. Read together under the Lectionary's pairing, the
two passages show one sovereignty: the Lord who took Cyrus by the hand, though
Cyrus did not know him, is the God in whose world Caesar's coin circulates;
tribute is therefore a real debt, and the same sovereignty bounds it.

\subsection{A gentile king called the Lord's anointed}

The first reading's opening words raise the difficulty. The Lord speaks ``to
my anointed Cyrus'', and twice in three verses he adds, ``thou hast not known
me''. The title that Israel gave its kings is given to a Persian who worships
other gods. How can a man be the Lord's anointed and not know the Lord?

Jerome answers from what the title meant. \latin{Iste appellatus est
christus, id est, unctus Domini, quod erat insigne apud Hebraeos regiae
potestatis, ut quomodo apud nos diadema et purpura solis imperatoribus datur:
sic apud Hebraeos regnaturi perfundebantur unguento}: ``this man is called
\emph{christus}, that is, anointed of the Lord, which was the mark of royal
power among the Hebrews; as among us the diadem and the purple are given to
emperors alone, so among the Hebrews those who were to reign were poured over
with ointment'' (\work{Commentarii in Isaiam} XII, PL~24, col.~457). Saul too
was called the Lord's anointed. The title names an office, not a faith. Aquinas
says the same in fewer words: \latin{Christo, idest regi \ldots\ propter
regiam dignitatem}, ``to Christ, that is, to the king \ldots\ because of his
royal dignity'' (\work{Expositio super Isaiam} c.~45, Parma ed., p.~539).

Both then ask what the Lord wanted of this king's power. Jerome reads the
omitted v.~3, the hidden treasures, as given ``that you who formerly
worshipped idols might perceive the one God through his benefits''
(\latin{ut qui prius colebas idola, beneficiis unum sentires Deum}); ``thou
hast not known me'' means \latin{simulacra coluisti, non Deum}, ``you
worshipped images, not God''. He reports Josephus's account that Cyrus read
the prophecy written about him, and on v.~6 he says that Cyrus wrote to all
the nations that there is no God but the God of Israel, or else that the
captives' release made God known to all (cols.~457--458). Aquinas states the
purpose as a fruit the Lord required: the fruit of the gift of power in Cyrus
himself was the knowledge of God, \latin{suae cognitionis}; the fruit in
others, that all from east to west should know there is no God besides him,
\latin{Quod completum fuit in hoc quod Cyrus epistolis suis laudem Dei scripsit
ubique; vel quia omnes Deum collaudaverunt videntes Iudaeorum liberationem},
``which was fulfilled in that Cyrus wrote the praise of God everywhere in his
letters, or because all praised God seeing the liberation of the Jews''
(p.~539). John Chrysostom, explaining how the Magi were moved to seek the
child, adds a qualification about how God moves such a king. God moved their
soul, ``which same kind of thing He did also in the case of Cyrus, disposing
him to let the Jews go. He did not however so do this as to destroy their free
will'' (\work{Homilies on Matthew}~6).

What the first reading shows, on these readings, is a God who uses the power of
those who do not know him. He grasps the king's hand, gives him victories and
treasures, and does it for the sake of Jacob his servant and so that the
nations may know that there is no other God. The king's power is real and given,
and its purpose is not the king's.

\subsection{``To Christ my Lord'': a reading Jerome rejects}

The verse has a second history. Some early Latin and Greek readers took it as
spoken to the Son. Cyprian's collection of testimonies cites it, among the
proofs that the gentiles would believe in Christ, as ``Thus saith the Lord God
to Christ my Lord, whose right hand I hold, that the nations may hear Him''
(\work{Ad Quirinum} I.21); Novatian cites ``Thus saith the Lord to Christ my
Lord'' among proofs that the Son is a second person (\work{De Trinitate} 26);
and the \work{Epistle of Barnabas} reads ``The Lord said to Christ, my Lord,
whose right hand I have holden'' (12). The editor of Novatian's English notes
the source of the reading: a copyist wrote the Greek word for ``Lord'',
\textit{Kyri\=o}, where the verse had the name of Cyrus, \textit{Kyr\=o}.

Jerome names the error and refuses it. \latin{Scio in hoc capitulo non solum
Latinorum, sed Graecorum plurimos vehementer errare, existimantium scriptum
esse: Sic dicit Dominus Christo meo, Domino}: ``I know that in this chapter
very many, not only of the Latins but of the Greeks, err greatly, thinking it
written, Thus saith the Lord to my Christ, the Lord''. It is said not
\textit{Kyri\=o} but to Cyrus, king of the Persians. And he will not let ``thou
hast not known me'' be said of Christ: \latin{In quo loco satis mirari nequeo,
quae stultitia sit legentium, ut haec ad Christum referant, per quem mundus
reconciliatus est Deo}, ``I cannot wonder enough at the folly of readers who
refer these things to Christ, through whom the world was reconciled to God''
(col.~457). The Lectionary's Latin, \latin{de uncto meo Cyro}, names Cyrus and
excludes the variant.

A narrower difference remains. Jerome, having refused the variant, still reads
the Vulgate's \latin{assimilavi te} of Cyrus as going before ``as a type and
image'' of the Christ to whom he was likened. Aquinas reads the same word of
the other saints whose names the Lord foretold (\latin{assimilavi, hoc in aliis
sanctis meis, quorum nomina prius a Domino praedicta sunt}) and sees no type.
The Lectionary's English, following the Hebrew, has no likeness at all. On
the Lectionary's words Cyrus is the Lord's instrument; the figure of Christ is
Jerome's, drawn from the Vulgate's word.

\subsection{Render, not give}

If the Lord's sovereignty reaches a pagan king, what follows for the coin in
the questioners' hands? Chrysostom answers from the verb. The questioners had
asked whether it was lawful to \emph{give} tribute; Jesus told them to
\emph{render} it. ``For this is not to give but to render, and this He shows
both by the image, and by the superscription'' (\work{Homilies on Matthew}
70.2). The coin bears Caesar's image and name; it is already his, and paying it
is returning what is owed, not making a gift. Chrysostom then reads the second
clause as protecting the first from misunderstanding: ``Then that they might
not say, Thou art subjecting us to men, He added, `And unto God the things
that are God's.' For it is possible both to fulfill to men their claims and to
give unto God the things that are due to God from us.'' He cites Paul, from a
chapter this Mass does not read: ``Render therefore to all men their dues. Tribute, to whom tribute is due''
(Rom 13:7).

In his homily on that chapter of Romans Chrysostom makes the same point from
the other side:
Paul wrote ``to show that it was not for the subversion of the commonwealth
that Christ introduced His laws, but for the better ordering of it'', and
``this subjection is not subversive of religion'' (\work{Homilies on Romans}
23). Justin Martyr had told the emperors that
Christians pay ``the taxes both ordinary and extraordinary, as we have been
taught by Him'', quoting this saying, and then: ``Whence to God alone we
render worship, but in other things we gladly serve you'' (\work{First
Apology}~17).

Aquinas gives the double belonging its shortest form. Explaining the answer,
he paraphrases it: \latin{quasi dicat: Vos estis Dei, et Caesaris, et habetis
in usu vestro et quae Dei sunt, et quae Caesaris}, ``as if he said: You are
God's, and Caesar's, and you have in your use both the things that are God's
and the things that are Caesar's''. What we have from the world binds us, he
says, to keep peace with the world, and he cites the first chapter of Baruch,
where the exiles are bidden to pray for the life of the king of Babylon
(\work{Super Matthaeum} c.~22, Venice ed., p.~284).

\subsection{Where the debt ends}

The same sovereignty that makes the debt real sets its limit. Chrysostom says
so in the next breath: ``But thou, when thou hearest, `Render unto Caesar the
things which are Caesar's,' know that He is speaking only of those things,
which are no detriment to godliness; since if it be any such thing as this,
such a thing is no longer Caesar's tribute, but the devil's'' (70.2). The
second clause of the saying does not merely add a second obligation; it
measures the first. Whatever Caesar claims that belongs to God is no longer
Caesar's to claim.

Augustine gives the limit a case. Expounding a psalm on the rod of the wicked
over the lot of the just, he says that a Christian servant should serve even a
bad master, and that what holds of master and slave holds ``also of powers and
kings''. Then: ``Julian was an infidel Emperor, an apostate, a wicked man, an
idolater; Christian soldiers served an infidel Emperor; when they came to the
cause of Christ, they acknowledged Him only who was in heaven. If he called
upon them at any time to worship idols, to offer incense; they preferred God
to him: but whenever he commanded them to deploy into line, to march against
this or that nation, they at once obeyed. They distinguished their everlasting
from their temporal master; and yet they were, for the sake of their
everlasting Master, submissive to their temporal master'' (\work{Enarrationes
in Psalmos} 124.7). The soldiers' obedience and their refusal come from the
same source: they obeyed the emperor for the sake of the master in heaven, and
for his sake refused the incense. Elsewhere Augustine asks what harm
Christians had done to the kingdoms of the earth: ``Did their King forbid His
soldiers to pay and to render due service to the kings of the earth?'', and he
answers with this saying and with the tribute paid from the fish's mouth
(118, under the letter Schin).

On this reading the saying is neither a charter for two independent realms nor
a refusal of civil duty. Caesar's claim is real because it lies within a world
where the Lord alone is God, and for the same reason it is bounded.

\subsection{A disagreement about the ground of the debt}

Not every witness grounds the debt in the same way. Hilary of Poitiers makes the obligation depend on holding what is
Caesar's: \latin{Si enim nihil eius penes nos resederit, conditione reddendi ei
quae sua sunt non tenebimur}, ``for if nothing of his remains with us, we shall
not be held by the condition of rendering to him what is his''. But if we rest
on his goods and use the right of his power, then it is beyond complaint that
Caesar be given back what is Caesar's (\work{Commentarius in Matthaeum} 23.2).
Ambrose, on Luke, draws the consequence for the disciple: \latin{Et tu si vis
non esse obnoxius Caesari, noli habere quae mundi sunt: sed si habes divitias,
obnoxius es Caesari. Si vis nihil regi debere terreno, relinque omnia tua, et
sequere Christum}, ``If you do not wish to be bound to Caesar, do not have the
things of the world; but if you have riches, you are bound to Caesar. If you
wish to owe nothing to an earthly king, leave all that is yours and follow
Christ'' (\work{Expositio evangelii secundum Lucam} IX.35).

For Hilary and Ambrose, then, the debt follows possession: whoever keeps the
world's goods owes the world's ruler, and the one who has left everything owes
him nothing. Chrysostom and Augustine say otherwise, each at a passage already
quoted. Chrysostom says that Paul opens his command with ``let every soul be
subject unto the higher powers'' ``to show that these regulations are for all,
even for priests, and monks, and not for men of secular occupations only'':
the subjection holds ``if thou be an Apostle even, or an Evangelist, or a
Prophet, or anything whatsoever'' (\work{Homilies on Romans} 23). The monk
and the Apostle, who have done what Ambrose bids the disciple do, are named
among those who owe. Augustine, in the exposition that gives the
soldiers under Julian, grounds the honour paid to rulers in God's ordering of
his Church: ``thus hath God ordained His Church, that every power ordained in
the world may have honour, and sometimes from those who are better than those
in power'' (\work{Enarrationes in Psalmos} 124.7). For them the debt does not
wait on what a Christian holds of Caesar's; it binds the one who has left
everything as well, and its limit is fixed not by renunciation but by what
belongs to God. The two positions agree that the tribute is lawful and that
God's claim comes first. They disagree about whether one who keeps nothing of
the world's goods still owes Caesar, and so about what the debt rests on.

\subsection{The Lord reigns among the nations}

The psalm turns the first reading's ``there is none besides me'' into a summons
to the nations. The Lord is to be feared above all gods, ``For all the gods of
the Gentiles are devils: but the Lord made the heavens'' (Ps 95:5). The Hebrew
calls them idols; the Greek and Latin, which Augustine and Bellarmine expound,
call them demons. Bellarmine: the gods of the gentiles ``are not true gods, but
demons''. Augustine reads the psalm's whole movement as the Redeemer's freeing
of the nations from captivity: ``All nations then had devils for their gods''
(\work{Enarrationes in Psalmos} 95.5).

``Say ye among the Gentiles, the Lord hath reigned'' (v.~10). Bellarmine
explains the kingship: ``God always reigns in heaven, and he reigns on earth
through his power and majesty; but he began to reign, through faith, among the
gentiles, from the coming of the Messias, where the devil previously reigned''.
Augustine's Old Latin psalter read the verse with an addition, ``the Lord hath
reigned from the wood'', which the appointed Latin does not have; on that text,
not the verse sung at this Mass, he says that ``Christ already possesseth
regions where the Roman empire hath never yet reached; what is as yet closed to
those who fight with the sword, is not closed to Him who fighteth with wood''
(95.2). For one Father who read the psalm within the Roman empire, the Lord's
reign did not wait on Caesar's. The strophe ends with the judgment, ``he will judge the people with justice'' (10c), and Bellarmine,
at the psalm's close, describes the Judge as one who will ``have no regard of
persons, and \ldots\ judge in all justice''. The questioners had praised Jesus
for regarding no man's person (Mt 22:16); Bellarmine gives the same quality to
the Lord who judges the nations.

\subsection{The Missal's texts in this reading}

The Entrance Antiphon asks for shelter under the shadow of God's wings, which
Augustine reads as ``the defence of Thy love and mercy''. The protection asked
for is God's and not a ruler's. In the psalm the same verse opens by asking
protection from those who resist God's right hand, words the antiphon leaves
out; in the first reading it is the Lord who grasps the right hand of a
king.

The Collect asks to serve God's majesty, \latin{maiestati tuae \ldots\
servire}, with a sincere heart. The Latin word belonged to gods and rulers
alike: Lewis and Short give \latin{maiestas} first ``of the gods'', then of
men in high station and of the sovereignty of the Roman people, and also ``as
a title of honor of the Roman emperors'' (\work{Latin Dictionary}, s.v.
\latin{majestas}, p.~1101). In this interpretation the Collect's
word, said to God, stands beside the Gospel's Caesar: the majesty served
first is God's, and every other service, Caesar's included, is ordered under
it.

The second reading greets a church that had learned what it meant to live
under Caesar's decrees and confess another king. Acts, which is not read at
this Mass, tells how Jason and certain brethren of Thessalonica were dragged
before the city's rulers with the charge that ``these all do contrary to the
decrees of Caesar, saying that there is another king, Jesus'' (Acts 17:7).
Chrysostom, expounding the church's ``labour of charity'', recalls that scene: ``Did not
they that warred against the Preaching \ldots\ drag Jason before the rulers of
the city?'' Acts reports the charge as the cry of those who raised the
tumult; whatever its intent, the church's faith and charity had to endure the
suspicion that it owed Caesar nothing.

The Alleluia verse places the faithful in the midst of a crooked and perverse
generation, shining as lights. Chrysostom's stars, which shine the brighter in
the night (\work{Homilies on Philippians}~8), describe a church living among
powers it does not control.

The Prayer over the Offerings asks to serve God's gifts with a free mind. In
this interpretation the freedom is that of those whose first service is God's:
they pay Caesar what is his without belonging to him.

The first Communion antiphon takes the eyes of the Lord upon those who fear
him from a psalm whose preceding verse says, ``The king is not saved by a great
army'' (Ps 32:16), a verse not sung. Aquinas, commenting on it, says that no temporal power can bring men to the salvation of the just
(\latin{nulla potestas temporalis potest homines ad salutem iustorum
perducere}), and cites ``Put not your trust in princes'' (Ps 145:2;
\work{Postilla super Psalmos} 32). The salvation the antiphon promises comes
from the Lord's regard, not from the strength of kings.

The second Communion antiphon, where it is sung, comes from a passage that
sets the Son of Man against the rulers of the gentiles: ``they who seem to rule
over the Gentiles lord it over them \ldots\ But it is not so among you'' (Mk
10:42--43). Bede reads the passage as teaching that one comes to the height of
the virtues ``not by power, but by humility'' (\latin{cum ad summitatem
virtutum non potentia, sed humilitate veniatur}; \work{In Marci evangelium
expositio} III, PL~92, col.~236). Chrysostom, on Matthew's form of the saying,
gives the Son of Man's own reason: ``For being King of the powers above, I was
willing to become man \ldots\ even my life did I give a ransom; and for whom?
For enemies'' (\work{Homilies on Matthew} 65.4). The antiphon itself keeps only
the ransom; the contrast with the rulers lies in its unsung context. Where it
is sung, that context makes the King who rules by service the measure of every
authority, Caesar's included.

The Prayer after Communion asks that the assembly be aided by temporal benefits
and instructed by eternal ones. That is the order of goods within which
Caesar's coin has its place: it helps, and it does not teach. Aquinas's ``You
are God's, and Caesar's'' names the double belonging that the prayer puts in
order.

\subsection{Where this interpretation strains}

The interpretation stands on the strongest relation the books declare, and its
difficulty lies in that relation. The Fathers speak to each half separately:
Jerome and Aquinas to the Lord's hand on Cyrus, Chrysostom and Augustine to
the coin and its limit. The two halves meet in the \work{Ordo}'s titles, and
the sovereignty that joins them is this reading's inference from the pairing.
The Fathers' cases, moreover, are emperors and kings, tribute, incense and the
deployment of soldiers; what the soldiers under Julian teach is the order of
two masters, not a form of government, and the reading reaches no further.

\subsection{The four senses of this interpretation}

\begin{fourSenses}
\item[Literal.] The Lord raises Cyrus, a king who does not know him, to subdue
nations and free Israel, so that all may know there is no other God (Is 45:1,
4--6); Jesus declares the census coin owed to Caesar and God's due owed to God
(Mt 22:21). Tribute is a debt rendered and not a gift (Chrysostom).
\item[Allegorical.] Christ is the King who reigns over the nations, by faith
among the gentiles (Bellarmine on Ps 95) and, in Augustine's Old Latin
psalter, ``from the wood''; where the second Communion antiphon is sung, he is
the Son of Man who gives his life as a ransom, whom the antiphon's unsung
context sets against the rulers of the gentiles (Bede; Chrysostom on the
parallel). Cyrus figures Christ only in Jerome, on the Vulgate's
\latin{assimilavi}.
\item[Moral.] Render civil dues as justice requires and keep peace with the
world from which you hold its goods (Aquinas), but refuse whatever is against
God: the soldiers who served Julian and would not offer incense (Augustine);
the limit at godliness (Chrysostom).
\item[Anagogical.] The Lord ``will judge the people with justice'' (Ps
95:10c), without regard of persons (Bellarmine), and the rod of the wicked over
the just is felt ``for a season'' and ``will not be there for ever''
(Augustine, \work{Enarrationes in Psalmos} 124.8).
\end{fourSenses}
